\documentclass[11pt]{article}

% ---------- Packages ----------
\usepackage[margin=1.9cm]{geometry}
\usepackage[T1]{fontenc}
\usepackage[utf8]{inputenc}
\usepackage{times}
\usepackage{xcolor}
\usepackage{enumitem}
\usepackage{tcolorbox}
\usepackage{fontawesome5}
\usepackage{hyperref}

% ---------- Palette ----------
\definecolor{accent}{HTML}{1D4E89}   % deep academic blue
\definecolor{accent2}{HTML}{2E9E83}  % teal for the AI box
\definecolor{todored}{HTML}{C0392B}

\newcommand{\todo}[1]{\textcolor{todored}{\textit{[#1]}}}
\newcommand{\contribhead}[2]{%
  {\bfseries\color{accent}\large #1\ #2}\\[1pt]
  {\color{accent}\rule{\linewidth}{1.1pt}}\\[4pt]%
}

\setlist[itemize]{leftmargin=1.2em, itemsep=1pt, topsep=1pt, partopsep=0pt}
\setlength{\parskip}{0pt}
\setlength{\parindent}{0pt}

\hypersetup{colorlinks=true, linkcolor=accent, urlcolor=accent, pdftitle={Contribution Statement -- ChartPRM}}

\pagestyle{empty}

\begin{document}

% ---------- Header ----------
\begin{center}
  {\LARGE\bfseries\color{accent} Contribution Statement}\\[0.25em]
  {\large ChartPRM: Process Reward Modeling and Step-Level Preference Optimization\\ for Complex Multimodal Chart Reasoning}\\[0.4em]
  {\small\itshape Process Reward Modelling Seminar \textbullet\ University of Heidelberg \textbullet\ \today}
\end{center}

\vspace{0.3em}
{\color{accent}\hrule height 1.4pt}
\vspace{0.6em}

% ---------- Team block ----------
\noindent
\begin{tabular}{@{}p{0.48\textwidth}p{0.48\textwidth}@{}}
  {\bfseries\faUser\ Yahor Lahunovich} & {\bfseries\faUser\ Ertugrul Taparci} \\[2pt]
  \small University of Heidelberg, Warsaw Univ.\ of Technology & \small University of Heidelberg \\
  \small Matriculation No.\ \todo{FILL IN} & \small Matriculation No.\ 4756036 \\
\end{tabular}

\vspace{0.5em}
\noindent\textbf{\faInfoCircle\ Both team members contributed meaningfully to this project.} Responsibilities and completed tasks are listed below.

\vspace{0.6em}

% ---------- Two-column contributions ----------
\noindent
\begin{minipage}[t]{0.485\textwidth}
\contribhead{Yahor Lahunovich}{}
\small
Built the core ChartPRM pipeline on \textbf{CharXiv} and authored the majority of the report.
\begin{itemize}
  \item Curated and stratified the 500-question CharXiv reasoning pool (8 domains, 62 chart types) and the 100-question holdout.
  \item Built the structured chain-of-thought rollout generation protocol for \texttt{Qwen2.5-VL-3B-Instruct}.
  \item Integrated the \texttt{muse-spark-1.1} multimodal LLM-as-judge with rollout-level batching.
  \item Built the preference-dataset curation and training pipelines for SFT, Full DPO, Step-DPO, KTO, and SFT$\rightarrow$DPO.
  \item Ran and evaluated all five systems on the CharXiv holdout.
  \item Built the visualization and figure-generation system used throughout the report.
  \item Co-authored and revised the project report throughout, including writing, structure, and the bibliography.
\end{itemize}
\vspace{2pt}
\todo{Yahor -- please revise this to your own words and add your own AI-assistance disclosure, if any (\texttt{.cursorrules}/\texttt{.agents/AGENTS.md} list Cursor, Gemini, Claude, and Antigravity as intended tools from early on).}
\end{minipage}%
\hfill
\begin{minipage}[t]{0.485\textwidth}
\contribhead{Ertugrul Taparci}{}
\small
Presented Yin et al.'s \textit{Dynamic and Generalizable Process Reward Modeling} (DG-PRM) for the seminar and directed its adaptation, plus a cross-dataset study.
\begin{itemize}
  \item Directed the reward-tree construction: 33 scored criteria under 9 categories, from clustered judge critiques.
  \item Directed the dynamic multi-criteria judge (scored via \texttt{gemini-3.5-flash-lite}, distinct from the \texttt{muse-spark-1.1} binary judge) and Pareto-dominance-filtered pair selection for DPO (\textbf{Pareto-DPO}).
  \item Directed the reference-free \textbf{SimPO} baseline implementation.
  \item Directed the ChartGemma specialist-baseline and oracle-ceiling diagnostic experiments.
  \item Designed and ran the cross-dataset generalization test on \textbf{ChartQA}: sampling, rollout generation, unmodified reward-tree transfer, and retraining/evaluating all seven methods on a 30-question holdout.
  \item Authored the DynamicPRM, Test-Time Verification and Search, and Cross-Dataset Generalization (ChartQA) sections, and their tables/figures; also revised the Abstract, Introduction, and ChartPRM Framework as the project's scope grew to two benchmarks and seven methods.
  \item \textbf{AI usage summary:} Used \textbf{Claude Code} (Anthropic) as an AI coding agent for essentially all of the code and analysis above (reward tree, dynamic judge, SimPO, ChartQA pipeline) and for drafting my report sections. I directed every experiment (what to try, when to pivot), set the evidentiary bar (e.g., requiring bootstrap significance tests before claiming a win), and personally verified and interpreted all results before they went into the report.
\end{itemize}
\end{minipage}

\vspace{1.1em}
\noindent{\small\itshape Both authors confirm this statement accurately reflects their respective contributions to the ChartPRM project.}

\end{document}
